\section{Constraints and Limitations}
\label{sec:properties}

The properties of large language models stated in this section constrain what any system built around them can guarantee. They are consequences of how the model is trained and queried, and they are stated here because the rest of this work depends on them.

\paragraph{Stochasticity.} Each output of a large language model is sampled from a distribution over tokens rather than read from a deterministic policy, according to the decoding strategy of Section~\ref{subsec:decoding}. Two runs on the same input can therefore return different outputs. Song et al. showed that this variability is large enough to affect benchmark conclusions \cite{songGoodBadGreedy2024}. The same model can score differently under greedy decoding and under sampling, and the gap is observed across model families and benchmarks \cite{songGoodBadGreedy2024}. The same analysis reported that the ranking of two models can change with the decoding configuration, so a comparison between models holds only for a fixed configuration \cite{songGoodBadGreedy2024}. The effect compounds in a system that chains many model calls: the variance of each call carries into the next, so two runs on the same task can reach different outcomes through different paths.

\paragraph{Context dependence.} Model behavior depends on the prompt and on whatever else occupies the context window, not on encoded rules or a stable policy. The position of the relevant content in the context is one factor. By keeping the input fixed and moving the document that contains the answer, Liu et al. reported that performance is highest when that document is at the beginning or at the end of the input context and degrades when the model must use it in the middle \cite{liuLostMiddleHow2023a}. The performance of GPT-3.5-Turbo on multi-document question answering drops by more than 20\% under this change, and when the answer document is placed in the middle of the context (the worst position), performance is lower than that of the same model with no documents supplied at all \cite{liuLostMiddleHow2023a}. Models with an extended context window are not better at using their input context than their non-extended counterparts \cite{liuLostMiddleHow2023a}. The effective policy is implicit in the model and in the context assembled around it, and it shifts whenever any part of that context changes.

\paragraph{Instruction following.} Instructions in the prompt state what the model is asked to do, and those instructions are not always followed. Zhou et al. measured this on instructions whose satisfaction a program can check, such as a word limit or a required keyword, using 25 instruction types over about 500 prompts \cite{zhouInstructionFollowingEvaluationLarge2023}. GPT-4 satisfies every verifiable instruction of a prompt in 76.89\% of the prompts, and PaLM 2 S in 43.07\% \cite{zhouInstructionFollowingEvaluationLarge2023}. The instructions measured this way are the ones that are easiest to state without ambiguity. The same work notes that language is subjective, and that the standard for an instruction such as writing with a funny tone is unclear \cite{zhouInstructionFollowingEvaluationLarge2023}.

\paragraph{Hallucination.} Large language models generate content that conflicts with the source available to them, called intrinsic hallucination, or that cannot be verified against that source, called extrinsic hallucination \cite{zhaoSurveyLargeLanguage2026, huangSurveyHallucinationLarge2025}. Hallucination occurs widely in existing models, including the most capable ones \cite{zhaoSurveyLargeLanguage2026}, and the false content is produced with high confidence \cite{wangSurveyLargeLanguage2024}.
